\begin{abstract}

\begin{center}
    \large{% Единое место для указания названия работы
Реализация векторизации ациклических участков кода по методу Superword-Level Parallelism в компиляторе языка программирования Go
} \\
    \large\textit{Самойлов Арсений Сергеевич} \\[1 cm]

    В работе рассматривается задача автоматической векторизации программного кода в компиляторе языка Go. 

    Отсутствие в стандартном компиляторе оптимизаций, использующих SIMD-инструкции, 
    ограничивает производительность приложений. 
    В работе предложен и реализован проход SLP-векторизации, адаптированный к SSA представлению и 
    архитектурным ограничениям Go: явному моделированию памяти, отсутствию сложных анализов и требованию быстрой компиляции.
    Разработанный векторизатор выполняет поиск изоморфных независимых операций, расширение групп, валидацию и 
    генерацию 128-битных векторных инструкций для архитектур x86-64 (SSE/AVX) ARM64 (NEON).
    Экспериментальная оценка на синтетических тестах и реальных приложениях 
    показала ускорение до XX\% при незначительном увеличении времени компиляции (порядка 1\%). 
    Работа демонстрирует принципиальную возможность и практическую пользу внедрения SLP-векторизации в основной компилятор Go.
    \vfill
\end{center}

\end{abstract}